% =============================================================================
% Pillar 3 — measured group-size evidence.
% Driver: platform_modal/scripts/group_size_review_1ed8_audit.py
% Inputs:  experiments/results/groupsize_zvf_sweep.json (G in {2,4,8,16}, n=3)
%          experiments/results/zvf_iter46_per_prompt_isog.tsv (505-task cohort)
% The earlier 20-cell Qwen3-8B/GSM8K "token-budget" grid (FALLBACK_ROWS in
% experiments/group_size_token_normalized.py) was hand-entered with no
% measured runs behind it; it and every analysis built on it are removed.
% =============================================================================
\section{Group Size: Measured Near-Ceiling Sweep}
\label{sec:group-size}

\paragraph{Question.}
Is GRPO's group size $G$ a consequential hyperparameter? Wu
et~al.~\cite{wu2025grpodpo} show that the group-centred GRPO update is an implicit
contrastive (DPO-like) objective in which $G$ mainly affects Monte-Carlo variance,
and report that 2-GRPO retains $97.6\%$ of 16-GRPO performance at $12.5\%$ of the
rollouts. The opposite intuition---that larger $G$ amortises wasted (all-correct or
all-incorrect) groups---is also common. This section reports a measured sweep on a
saturated task; Section~\ref{sec:group-size-scaleup} adds an unsaturated,
iso-compute sweep that includes $G{=}32$ and a Qwen3-8B/GSM8K check.

\paragraph{Sweep protocol.}
The sweep (\texttt{experiments/modal/modal\_groupsize\_zvf\_sweep.py}, one NVIDIA
A10 per run) trains Qwen2.5-0.5B on two-operand addition (operands uniform in
$\{1,\dots,99\}$) with LoRA (rank $16$, $\alpha{=}32$ on $q$/$v$), Adam at
learning rate $10^{-4}$, $16$ fresh prompts per optimizer step and $G$ completions
per prompt (temperature $1.0$, top-$p$ $1.0$, at most $10$ new tokens), for $40$
steps, at seeds $\{42,123,456\}$. Advantages are
$\hat A_i=(r_i-\bar r)/(\sigma_r+10^{-6})$ with the population standard deviation,
no KL penalty and no ratio clipping. The step count is fixed, so $G{=}16$ uses
$8\times$ the rollouts of $G{=}2$. Held-out accuracy is greedy on $200$ freshly
sampled problems; the last-10 reward is the mean training reward over the final
ten steps.

\paragraph{Held-out accuracy is flat in $G$.}
Table~\ref{tab:groupsize-summary} summarises the sweep. Held-out accuracy is
$0.9817$, $0.9883$, $0.9900$ and $0.9783$ for $G{=}2,4,8,16$; the numerical
maximum at $G{=}8$ is not distinguishable from its neighbours at $n{=}3$ seeds
(overlapping $1.96\cdot\mathrm{SE}$ intervals). $G{=}2$ retains $100.34\%$ of
$G{=}16$ held-out accuracy (difference $+0.0033$, bootstrap 95\% CI
$[-0.0083, 0.0150]$, $B{=}4000$; Table~\ref{tab:groupsize-dpo-check}), above the
$97.6\%$ Wu et~al.\ headline. A held-out TOST at $\varepsilon{=}0.02$ gives
$p{=}0.100$, so equivalence is not certified; with $n{=}3$ the minimum detectable
effect is $d{=}3.07$. Last-10 training reward is likewise flat
($0.971$--$0.984$). This is a saturated regime: near the ceiling, variance is
compressed and ``no detectable effect'' is the most the data can say.

\begin{table}[htbp]
\centering
\small
\begin{tabular}{lcccc}
\toprule
Metric & $G{=}2$ & $G{=}4$ & $G{=}8$ & $G{=}16$ \\
\midrule
Held-out acc.\ (mean $\pm$ SE, $n{=}3$) & $0.9817{\pm}.004$ & $0.9883{\pm}.002$ & $0.9900{\pm}.003$ & $0.9783{\pm}.006$ \\
Last-10 training reward & $0.9771$ & $0.9714$ & $0.9732$ & $0.9837$ \\
Mean ZVF & $0.838$ & $0.764$ & $0.691$ & $0.631$ \\
Rollouts per run & $1{,}280$ & $2{,}560$ & $5{,}120$ & $10{,}240$ \\
\bottomrule
\end{tabular}
\caption{Measured group-size sweep, Qwen2.5-0.5B / two-operand addition, 40
steps, 3 seeds per $G$
(\texttt{experiments/results/groupsize\_zvf\_sweep.json}). Held-out accuracy and
training reward are flat in $G$; the zero-variance fraction falls monotonically.}
\label{tab:groupsize-summary}
\end{table}

\begin{table}[htbp]
\centering
\small
\begin{tabular}{lcccc}
\toprule
Comparison & Mean A & Mean B & Diff. (A--B) & 95\% CI on diff. \\
\midrule
$G{=}2$ vs $G{=}16$ (Wu headline) & $0.9817$ & $0.9783$ & $+0.0033$ & $[-0.0083, +0.0150]$ \\
$G{=}4$ vs $G{=}16$               & $0.9883$ & $0.9783$ & $+0.0100$ & n/a \\
$G{=}8$ vs $G{=}16$               & $0.9900$ & $0.9783$ & $+0.0117$ & n/a \\
\bottomrule
\end{tabular}
\caption{Held-out comparisons on the measured sweep. Means to four decimals;
differences from unrounded means. $G{=}32$ on an unsaturated task is reported in
Section~\ref{sec:group-size-scaleup}.}
\label{tab:groupsize-dpo-check}
\end{table}

\paragraph{ZVF: measured vs closed form.}
Table~\ref{tab:groupsize-zvf-theory} compares the measured mean ZVF with the
homogeneous closed form $\mathrm{ZVF}(p,G)=p^G+(1-p)^G$, evaluated at each step's
batch training pass rate $p_t$ and averaged over steps and seeds, so both columns
describe the same rollouts. The measured ZVF exceeds the closed form at every $G$
(residual $+0.010$ to $+0.109$). This is the sign expected from prompt-difficulty
heterogeneity: by strict convexity of $p\mapsto p^G+(1-p)^G$,
$\mathbb{E}_x[p_x^G+(1-p_x)^G] > \bar p^{\,G}+(1-\bar p)^G$, so collapsing per-prompt
pass rates to a batch mean under-predicts ZVF. Conditional on a prompt and the
policy, i.i.d.\ decoding gives zero within-group correlation; the excess is a
marginal effect of mixing prompts of different difficulty. On the unsaturated
task of Section~\ref{sec:group-size-scaleup} the same gap is an order of magnitude
larger.

\begin{table}[htbp]
\centering
\small
\begin{tabular}{cccccc}
\toprule
$G$ & Mean training $p_t$ & ZVF (measured) & ZVF (closed form at $p_t$) & Residual & $1-\mathrm{ZVF}$ \\
\midrule
2  & $0.840$ & $0.838$ & $0.828$ & $+0.010$ & $0.162$ \\
4  & $0.863$ & $0.764$ & $0.709$ & $+0.055$ & $0.237$ \\
8  & $0.869$ & $0.691$ & $0.597$ & $+0.094$ & $0.309$ \\
16 & $0.873$ & $0.631$ & $0.522$ & $+0.109$ & $0.369$ \\
\bottomrule
\end{tabular}
\caption{Measured ZVF vs.\ the homogeneous closed form on the Qwen2.5-0.5B
addition sweep. Source:
\texttt{experiments/results/group\_size\_review\_1ed8\_zvf\_theory.tsv}.}
\label{tab:groupsize-zvf-theory}
\end{table}

\paragraph{Matched-rollout trajectory check on the Wu pairing.}
A two-seed panel on exactly $G{=}2$ vs.\ $G{=}16$ at matched rollouts (2{,}560 per
arm: $G{=}2\times160$ steps vs.\ $G{=}16\times20$ steps; Qwen3-8B/GSM8K via Tinker,
LoRA rank~4, batch~8, 512-token completions, learning rate $3\times10^{-5}$ in both
arms, seeds 123/456; \texttt{tinker-runs/results/er2b\_*.json}) shows the $G{=}2$
arms driving training reward to $\approx0.9$--$1.0$ and ending in an all-correct,
zero-variance regime ($\mathrm{ZVF}\approx0.75$--$1.0$), while the $G{=}16$ arms end
mid-learning (training reward $\approx0.3$--$0.5$, $\mathrm{ZVF}\le0.25$). Group size
is confounded here with an $8\times$ difference in optimizer steps at a fixed
learning rate, and held-out accuracy was not evaluated for these arms, so the panel
describes training dynamics only.

\paragraph{Conditional utility audit (505 prompt--seed tasks).}
\label{sec:p3-utility-audit}
A separate, purely analytical audit
(\texttt{analysis/breakthroughs\_2026-07-13/analyze\_breakthroughs.py}) uses the
$505$ distinct (problem, seed) pairs from the three base-model Qwen3-8B GSM8K
evaluations ($G{=}8$, no training; the same groups as the companion ZVF paper)
whose per-prompt pass rate lies in $(0.05, 0.95)$, i.e.\ every mixed group
(\texttt{experiments/results/zvf\_iter46\_per\_prompt\_isog.tsv}; mean
$p_x{=}0.674$). For each $G$ it computes the cohort-mean non-degenerate yield
$Y_G=\mathbb{E}_x[1-p_x^G-(1-p_x)^G]=1-\overline{\mathrm{ZVF}}$ under the
independent-completion model and scores it by $U(G)=Y_G/\sqrt{G}$. The $\sqrt{G}$
divisor is a \emph{declared} diminishing-returns proxy, not the generation cost,
which is linear in $G$; dividing by $G$ instead makes $G{=}2$ optimal
(Section~\ref{sec:frontier-group-size}). Under the $\sqrt{G}$ proxy,
$U(4){=}0.32914$ and $U(5){=}0.32737$, so $G{=}4$ is selected over $G{=}5$ by
$0.00177$ (task-resampling bootstrap 95\% interval $[0.00082, 0.00270]$;
re-verified by \texttt{tools/review\_ledger/P3\_group\_size\_audit.py}). Here $p_x$
is the plug-in $k/8$ from the same eight rollouts, so $Y_G$ inherits plug-in bias
and, for $G\neq8$, extrapolates from an eight-sample estimate; the all-correct and
all-wrong groups are excluded by construction.
Stratifying the same cohort by pass rate gives
$Y_{32}-Y_4={+}0.29$ ($p_x<0.5$, $n{=}62$), ${+}0.15$ ($0.5\le p_x<0.75$, $n{=}157$) and
${+}0.45$ ($p_x\ge0.75$, $n{=}286$). No training run was performed for this audit:
it describes how often a group of a given size would contain contrast, not whether
that contrast improves held-out accuracy.

\paragraph{Conclusions.}
\begin{enumerate}
  \item On the saturated addition task, held-out accuracy and training reward do
        not depend detectably on $G\in\{2,\dots,16\}$; $G{=}2$ retains $100.3\%$ of
        $G{=}16$, consistent with the Wu et~al.\ reading but not certified as
        equivalent ($n{=}3$).
  \item Larger $G$ reliably lowers the zero-variance fraction, and the measured
        ZVF exceeds the homogeneous closed form, as prompt-difficulty
        heterogeneity predicts.
  \item Whether the lower ZVF of larger groups buys held-out accuracy on an
        unsaturated task at equal compute is tested in
        Section~\ref{sec:group-size-scaleup}.
\end{enumerate}
